\section{模組說明}
Portfolio 模組提供僅供研究的工作流：發現 \texttt{Model Portfolio}、維護 \texttt{My Portfolio}、生成 \texttt{AI Portfolio Recommendation}、分析 portfolio 草稿，以及根據歷史 \texttt{Baseline Snapshot} 審閱已保存的 \texttt{My Portfolio}。
模組由三個入口標簽組織：\texttt{Recommended} 用於瀏覽管理員維護的 Model Portfolio 並開始 recommendation quiz；\texttt{My portfolios} 用於創建、查看、編輯、刪除、分析和審閱用戶保存的 portfolio；\texttt{Watchlist} 用於查看與移除用戶保存供日後參考的 Model Portfolio。
\begin{quote}\small 範圍說明：本模組中的預測、分析、Recommendation 和 \texttt{Suggested actions} 僅供研究，不代表收益保證，也不會觸發券商訂單或交易執行。\end{quote}
\section{User Story /\allowbreak{} Video 索引}
\begingroup
\small
\begin{longtable}{L{0.18\linewidth}L{0.24\linewidth}L{0.28\linewidth}L{0.16\linewidth}}
\toprule
\textbf{Story} & \textbf{視頻主題} & \textbf{核心路徑} & \textbf{Video} \\
\midrule
\hyperlink{us-port-01}{US-PORT-01} & 探索並保存 Model Portfolio & \texttt{Recommended} -\textgreater{} \texttt{Model Portfolio} -\textgreater{} \texttt{Watchlist} /\allowbreak{} \texttt{My Portfolio} & \href{https://jjpvro70sief.jp.larksuite.com/wiki/SjhpwcDNLirUxpkHh1vjUbgqphY}{視頻} \\
\hyperlink{us-port-02}{US-PORT-02} & 創建、編輯和維護 My Portfolio & \texttt{My portfolios} -\textgreater{} \texttt{Edit} /\allowbreak{} \texttt{Baseline snapshots} & \href{https://jjpvro70sief.jp.larksuite.com/wiki/I9gxw6Uugiw5RxkSJ91jVHKipMf}{視頻} \\
\hyperlink{us-port-03}{US-PORT-03} & 生成 AI Portfolio Recommendation & \texttt{Find Your Ideal Portfolio} -\textgreater{} \texttt{Portfolio recommendation} & \href{https://jjpvro70sief.jp.larksuite.com/wiki/NSZ1wCHloiD8QAkYSDsj9CNtpje}{視頻} \\
\hyperlink{us-port-04}{US-PORT-04} & 分析 Portfolio 並審閱 AI 建議 & \texttt{Draft Portfolio} -\textgreater{} \texttt{Analysis} -\textgreater{} \texttt{Suggestions} & \href{https://jjpvro70sief.jp.larksuite.com/wiki/Ey4Bw3xDxiPeETk3yBBjhxOHpdf}{視頻} \\
\hyperlink{us-port-05}{US-PORT-05} & 根據 Baseline Snapshot 審閱 & \texttt{Portfolio Review} -\textgreater{} \texttt{Baseline} -\textgreater{} \texttt{Analysis} -\textgreater{} \texttt{Suggestions} & \href{https://jjpvro70sief.jp.larksuite.com/wiki/H1ULwBLKeiaEIKkEbQDjhRiepue}{視頻} \\
\bottomrule
\end{longtable}
\endgroup
\section{統一術語}
\begingroup
\small
\begin{longtable}{L{0.24\linewidth}L{0.68\linewidth}}
\toprule
\textbf{UI 原文} & \textbf{文檔含義} \\
\midrule
\texttt{Model Portfolio} & 由管理員維護並發布的 portfolio template，不是用戶持久化配置 \\
\texttt{My Portfolio} /\allowbreak{} \texttt{My portfolios} & 用戶創建並保存的研究型 portfolio，以及其列表標簽 \\
\texttt{Watchlist} & 用戶保存供日後參考的 Model Portfolio 列表 \\
\texttt{Add to My Portfolio} & 將 Model Portfolio 或 Recommendation 轉為新的或合並到已有 My Portfolio \\
\texttt{Baseline Snapshot} & 某個 My Portfolio 在特定時間保存的持倉、數量、價值和權重參考配置 \\
\texttt{Baseline History} /\allowbreak{} \texttt{Baseline snapshots} & 管理 Baseline Snapshot 的歷史抽屜/\allowbreak{}頁面 \\
\texttt{Analysis Workspace} & 對 portfolio 工作草稿執行 Holdings、Analysis、Diagnosis、Suggestions 流程的工作區 \\
\texttt{Suggested actions} & 算法生成的量化調整預覽；是否可編輯/\allowbreak{}應用取決於來源 \\
\texttt{Portfolio Review} & 將當前 My Portfolio 與所選 Baseline Snapshot 比較的工作流 \\
\texttt{Recommendation Build} & 展示 Recommendation 的風險畫像、篩選、優化、候選範圍和警告 \\
\texttt{Data coverage} /\allowbreak{} \texttt{Unavailable} /\allowbreak{} \texttt{Partial} & 資料不完整或指標不可用時的明確邊界狀態 \\
\bottomrule
\end{longtable}
\endgroup
\prdhr
\hypertarget{us-port-01}{}
\section{User Story 1 - 探索並保存 Model Portfolio}
\textbf{Story ID:} \texttt{US-PORT-01}
\textbf{用戶故事：} 作為用戶，我希望檢查精選的 Model Portfolio，了解其配置和示意性預測，將其保存以便稍後參考，並使用 Add to My Portfolio。
\textbf{Video:} \href{https://jjpvro70sief.jp.larksuite.com/wiki/SjhpwcDNLirUxpkHh1vjUbgqphY}{US-PORT-01 驗收視頻}
\subsection{Walkthrough}
\begin{enumerate}
\item 打開 \texttt{Portfolio \textgreater{} Recommended}，瀏覽 Model Portfolio 卡片及其 \texttt{Reference} 說明。
\item 打開一個 Model Portfolio，查看回報、風險、期限、配置和 ETF allocations。
\item 瀏覽 \texttt{Projected Asset Growth}，操作時間範圍和歷史資料點，並查看可用的 backtest 與 data coverage。
\item 將 Model Portfolio 保存到 \texttt{Watchlist}，確認可以從 Watchlist 重新打開或移除。
\item 選擇 \texttt{Add to My Portfolio}，在 \texttt{Visual} 和 \texttt{ETF list} 模式中審閱或調整金額、權重和份額。
\item 使用 \texttt{Add as new} 創建新的 My Portfolio，並從 \texttt{My portfolios} 打開確認。
\item 使用 \texttt{Add to existing} 將配置加入已有 My Portfolio，並確認持倉結果更新。
\item 從 Model Portfolio 選擇 \texttt{Start analysis}，瀏覽只讀的分析與建議流程。
\end{enumerate}
\screenshotsTripleCompact{assets/screenshots/port-us1-01-model-portfolio-catalogue.png}{PORT-US1-01}{\texttt{Recommended} 目錄、recommendation quiz 入口、Model Portfolio cards、allocation mix 和 \texttt{Reference}。}{assets/screenshots/port-us1-02-model-portfolio-detail.png}{PORT-US1-02}{\texttt{Anchor Income} Model Portfolio 詳情，包含 return/\allowbreak{}risk/\allowbreak{}horizon、\texttt{Reference}、allocation map 和 projection 入口。}{assets/screenshots/port-us1-extra-etf-allocation.png}{Model Portfolio ETF allocation}{\texttt{Anchor Income} 的 \texttt{ETF allocation} 列表、ETF metadata、target weights 和 \texttt{Total 100\%}。}
\screenshotDescsThree{\textbf{PORT-US1-01} \textemdash{} \texttt{Recommended} 目錄、recommendation quiz 入口、Model Portfolio cards、allocation mix 和 \texttt{Reference}。}{\textbf{PORT-US1-02} \textemdash{} \texttt{Anchor Income} Model Portfolio 詳情，包含 return/\allowbreak{}risk/\allowbreak{}horizon、\texttt{Reference}、allocation map 和 projection 入口。}{\textbf{Model Portfolio ETF allocation} \textemdash{} \texttt{Anchor Income} 的 \texttt{ETF allocation} 列表、ETF metadata、target weights 和 \texttt{Total 100\%}。}
\screenshotsTripleCompact{assets/screenshots/port-us1-03-projection-backtest.png}{PORT-US1-03}{\texttt{Projected Asset Growth} 完成態，包含 history、P50、P25-P75、P2.5-P97.5、Today crosshair、summary 和 \texttt{Partial} data warning。}{assets/screenshots/port-us1-04-watchlist-saved.png}{PORT-US1-04}{\texttt{Watchlist} 中已保存的 \texttt{Anchor Income} Model Portfolio，包含 saved count、事實信息、allocation 和 \texttt{Reference}。}{assets/screenshots/port-us1-05-add-to-portfolio-editor.png}{PORT-US1-05}{\texttt{Add to My Portfolio} 的 Visual editor，包含 allocation overview、total portfolio value、sector weights 和兩個保存目標。}
\screenshotDescsThree{\textbf{PORT-US1-03} \textemdash{} \texttt{Projected Asset Growth} 完成態，包含 history、P50、P25-P75、P2.5-P97.5、Today crosshair、summary 和 \texttt{Partial} data warning。}{\textbf{PORT-US1-04} \textemdash{} \texttt{Watchlist} 中已保存的 \texttt{Anchor Income} Model Portfolio，包含 saved count、事實信息、allocation 和 \texttt{Reference}。}{\textbf{PORT-US1-05} \textemdash{} \texttt{Add to My Portfolio} 的 Visual editor，包含 allocation overview、total portfolio value、sector weights 和兩個保存目標。}
\screenshotsSingleCompact{assets/screenshots/port-us1-extra-list-editor.png}{Add to My Portfolio ETF list editor}{同一配置切換到 \texttt{ETF list} 後的 direct-share editor，展示 industry groups、ETF rows 和 share steppers。}
\screenshotDescSingle{\textbf{Add to My Portfolio ETF list editor} \textemdash{} 同一配置切換到 \texttt{ETF list} 後的 direct-share editor，展示 industry groups、ETF rows 和 share steppers。}
\subsection{Acceptance Criteria}
\begingroup
\small
\begin{longtable}{L{0.18\linewidth}L{0.30\linewidth}L{0.18\linewidth}L{0.12\linewidth}L{0.10\linewidth}}
\toprule
\textbf{ID} & \textbf{Requirement} & \textbf{Reference} & \textbf{Ownership} & \textbf{Priority} \\
\midrule
US1-AC01 & 可見 Portfolio 標題和三個標簽\textemdash{}\textemdash{}Recommended、My portfolios、Watchlist，且已選中 Recommended。 & US-PORT-01 & Integration & Must \\
US1-AC02 & 每張已載入的 Model Portfolio 卡片均使用後端資料展示標題、回報範圍、風險等級標簽、期限、資產類別數量和資產配置可視化。 & US-PORT-01 & Integration & Must \\
US1-AC03 & 每張 Model Portfolio 卡片顯示與該 Model Portfolio 對應的 Reference 說明。 & US-PORT-01 & Integration & Must \\
US1-AC04 & View details 打開所選 Model Portfolio，而非其他目錄項目。 & US-PORT-01 & Integration & Must \\
US1-AC05 & 首次打開或刷新 Model Portfolio 目錄時顯示載入狀態；完成後以當前返回的 Model Portfolio 替換占位。 & US-PORT-01 & Integration & Must \\
US1-AC06 & 詳情頭部展示所選 Model Portfolio 的標題、回報、風險、期限和摘要。 & US-PORT-01 & Integration & Must \\
US1-AC07 & 詳情事實值與目錄卡片及當前 Model Portfolio 資料一致，不沿用先前打開的 Model Portfolio 內容。 & US-PORT-01 & Integration & Must \\
US1-AC08 & 參考說明以专用的來源/\allowbreak{}參考樣式顯示，並與目錄卡片上的說明一致。 & US-PORT-01 & Integration & Must \\
US1-AC09 & 配置卡展示資產類別分組及其 ETF 代碼；顯示的 ETF 數量和資產類別數量與返回的詳情一致。 & US-PORT-01 & Integration & Must \\
US1-AC10 & ETF 配置區顯示 Total 100\%，每行展示代碼、名稱、資產類別標簽、目標權重，以及可用的發行方/\allowbreak{}最新價格/\allowbreak{}1 年回報元資料。 & US-PORT-01 & Integration & Must \\
US1-AC11 & 選擇 ETF 行會打開匹配的 ETF 詳情；返回僅關閉該浮層，不丟失 Model Portfolio 詳情狀態。 & US-PORT-01 & Integration & Must \\
US1-AC12 & 從 ETF 詳情返回後仍停留在當前 Model Portfolio 詳情，配置內容和已載入狀態保持不變。 & US-PORT-01 & Integration & Must \\
US1-AC13 & 載入狀態明確說明蒙特卡洛模擬正在運行；資料可用前，不得以編造的點顯示圖表。 & US-PORT-01 & Integration & Must \\
US1-AC14 & 圖表區分歷史、P50 預測、P25\textendash{}P75 和 P2.5\textendash{}P97.5 區間、Today、初始價值、軸標簽及預測標註。 & US-PORT-01 & Integration & Must \\
US1-AC15 & 圖表下方摘要值與圖表的起始價值、配置期限、P50 終點及 P2.5\textendash{}P97.5 範圍一致。 & US-PORT-01 & Integration & Must \\
US1-AC16 & 水平拖動會改變可見時間線，同時使視圖保持在可用的歷史/\allowbreak{}預測範圍內。 & US-PORT-01 & Integration & Must \\
US1-AC17 & 捏合會改變可見時間窗口；縮放至邊界後不會顯示可用日期範圍外的資料，軸標簽和繪制值不越出圖表區域。 & US-PORT-01 & Integration & Must \\
US1-AC18 & 選擇有效歷史點顯示對應時間和值的十字準線。僅歷史偏移點可提供回測操作，未來預測點不可提供。 & US-PORT-01 & Integration & Must \\
US1-AC19 & Generate backtest 顯示載入狀態，隨後添加含更新圖例的回測百分位區間/\allowbreak{}中位數和實際序列。 & US-PORT-01 & Integration & Must \\
US1-AC20 & Hide backtest 移除生成的回測，但不移除原始歷史或預測。 & US-PORT-01 & Integration & Must \\
US1-AC21 & 預測標記為 Partial 時顯示對應的 Data coverage，並繼續使用當前 Model Portfolio 的可用資料，不沿用其他 Model Portfolio 的圖表結果。 & US-PORT-01 & Integration & Must \\
US1-AC22 & 預測明確描述為示意性/\allowbreak{}僅供研究，不得表述為保證的預測。 & US-PORT-01 & Integration & Must \\
US1-AC23 & 保存請求成功後，詳情控件才從空心星標變為實心星標。 & US-PORT-01 & Integration & Must \\
US1-AC24 & Watchlist 計數增加，且 Model Portfolio 以相同標題、事實信息、配置和 Reference 出現在 Watchlist 中。 & US-PORT-01 & Integration & Must \\
US1-AC25 & 可從 Watchlist 打開已保存的 Model Portfolio。 & US-PORT-01 & Integration & Must \\
US1-AC26 & 移除已保存的 Model Portfolio 會使其從 Watchlist 消失並减少計數；當不再有已保存的 Model Portfolio 時，顯示明確的 No portfolios saved 狀態。 & US-PORT-01 & Integration & Must \\
US1-AC27 & 返回 Recommended 後再次打開同一 Model Portfolio，其星標狀態與 Watchlist 中的最終保存狀態一致。 & US-PORT-01 & Integration & Must \\
US1-AC28 & Add to My Portfolio 頁面為所選 Model Portfolio 打開，並展示配置概覽、Model Portfolio 標題、風險/\allowbreak{}回報背景及默認建模金額。 & US-PORT-01 & Integration & Must \\
US1-AC29 & 修改總價值會在保持當前歸一化配置權重的前提下，重新計算 ETF 金額和數量。 & US-PORT-01 & Integration & Must \\
US1-AC30 & 拖動行業邊界僅改變相鄰兩個分組；到達可拖動邊界後繼續拖動不再降低該分組，且各分組不出現負權重、總和保持 100\%。 & US-PORT-01 & Integration & Must \\
US1-AC31 & 拖動 ETF 拆分會改變該分組內 ETF，保持分組總和不變，並更新組內及總體配置百分比。 & US-PORT-01 & Integration & Must \\
US1-AC32 & 從 Visual 切換至 List 模式時，當前金額和權重會按界面顯示的最新價格轉換為 ETF 份額；修改份額會重新計算總價值、分組權重、ETF 權重及配置可視化。 & US-PORT-01 & Integration & Must \\
US1-AC33 & 從 List 切回 Visual 模式時，當前份額會轉換為對應的總金額和歸一化權重；反複切換不會丟失最新編輯，且兩種模式均不允許負份額或負權重。 & US-PORT-01 & Integration & Must \\
US1-AC34 & 進入確認步驟時，摘要使用最後一次 Visual/\allowbreak{}List 編輯後的金額、權重和份額，不恢複為初始配置。 & US-PORT-01 & Integration & Must \\
US1-AC35 & 確認步驟展示目標名稱、添加總額、ETF 數量、每只 ETF 的權重/\allowbreak{}金額/\allowbreak{}添加數量以及結果數量。 & US-PORT-01 & Integration & Must \\
US1-AC36 & 當 My Portfolio 名稱為空或創建請求進行中時，Confirm 被禁用。 & US-PORT-01 & Integration & Must \\
US1-AC37 & 創建成功後，新名稱出現在 My portfolios 中；持倉數與確認摘要一致，基礎貨幣與創建時選擇一致。 & US-PORT-01 & Integration & Must \\
US1-AC38 & 重新打開或刷新後，持倉、數量、配置和價值仍與確認結果一致。 & US-PORT-01 & Integration & Must \\
US1-AC39 & UI 說明已創建研究型 My Portfolio，未發生下單、交易或執行。 & US-PORT-01 & Integration & Must \\
US1-AC40 & 創建請求進行中時 Confirm 不可重複觸發；完成後列表只出現一個對應的新 My Portfolio。 & US-PORT-01 & Integration & Must \\
US1-AC41 & 目標選擇器列出當前 My Portfolio，並顯示各項的名稱、價值、持倉數和配置摘要。 & US-PORT-01 & Integration & Must \\
US1-AC42 & 預覽載入完成前 Confirm 不可用。 & US-PORT-01 & Integration & Must \\
US1-AC43 & 預覽展示已有價值、待添加 ETF 行、Model Portfolio 權重、金額、最新價格、添加數量和結果數量。 & US-PORT-01 & Integration & Must \\
US1-AC44 & 確認後更新所選已有 My Portfolio；列表和詳情刷新後顯示新的持倉結果。 & US-PORT-01 & Integration & Must \\
US1-AC45 & 重新打開目標後，預覽中的已有代碼顯示合並後的數量，預覽中的新增代碼顯示為新持倉。 & US-PORT-01 & Integration & Must \\
US1-AC46 & 從目標選擇或確認頁返回時，當前編輯配置保持不變；最終添加結果與最後確認的預覽一致。 & US-PORT-01 & Integration & Must \\
US1-AC47 & Analysis Workspace 在 Holdings 打開，來源標題與所選 Model Portfolio 一致，並按 Model Portfolio ETF 配置預載入代碼、份額、價值和歸一化權重。 & US-PORT-01 & Integration & Must \\
US1-AC48 & 對 Model Portfolio 草稿的編輯僅作用於工作副本，不會創建或修改 My Portfolio；由於來源不是 My Portfolio，建議結果不顯示 Apply to My Portfolio 控件。 & US-PORT-01 & Integration & Must \\
US1-AC49 & \texttt{Analyze draft portfolio} 使用當前編輯後的 Model Portfolio 草稿；分析、診斷和建議的共同行為符合用戶故事 4，且不會複用其他來源的舊結果。 & US-PORT-01 & Integration & Must \\
\bottomrule
\end{longtable}
\endgroup
\prdhr
\hypertarget{us-port-02}{}
\section{User Story 2 - 創建、編輯和維護 My Portfolio}
\textbf{Story ID:} \texttt{US-PORT-02}
\textbf{用戶故事：} 作為用戶，我希望創建和維護已保存的 ETF My Portfolio，包括其持倉和 Baseline Snapshot 歷史記錄。
\textbf{Video:} \href{https://jjpvro70sief.jp.larksuite.com/wiki/I9gxw6Uugiw5RxkSJ91jVHKipMf}{US-PORT-02 驗收視頻}
\subsection{Walkthrough}
\begin{enumerate}
\item 打開 \texttt{Portfolio \textgreater{} My portfolios}，查看已保存的 portfolios，並進入 \texttt{New Portfolio}。
\item 輸入名稱和 base currency，添加初始 ETF holdings，然後創建並打開新的 My Portfolio。
\item 查看 My Portfolio 的概覽、配置、預測和持倉，並保存一個 \texttt{Baseline Snapshot}。
\item 進入 \texttt{Edit}，調整持倉數量，並使用 \texttt{Add ETF} 或移除操作更新 portfolio。
\item 返回 \texttt{Baseline snapshots}，查看和管理編輯前後的歷史快照。
\item 從列表移除一個不再需要的 My Portfolio，並確認其他 portfolios 不受影響。
\end{enumerate}
\screenshotsTripleCompact{assets/screenshots/port-us2-01-my-portfolios-list.png}{PORT-US2-01}{\texttt{My portfolios} 列表、\texttt{New} 入口、saved portfolios、value、holdings count、asset mix 和 drift。}{assets/screenshots/port-us2-02-create-initial-holdings.png}{PORT-US2-02}{\texttt{New Portfolio} 初始狀態，包含 name、base currency、\texttt{Initial holdings}、Add ETF 入口和 research-only 邊界。}{assets/screenshots/port-us2-03-edit-holdings.png}{PORT-US2-03}{My Portfolio \texttt{Editing} 頁面，包含 value/\allowbreak{}holdings/\allowbreak{}status、weights、share steppers、remove controls 和 \texttt{Add ETF}。}
\screenshotDescsThree{\textbf{PORT-US2-01} \textemdash{} \texttt{My portfolios} 列表、\texttt{New} 入口、saved portfolios、value、holdings count、asset mix 和 drift。}{\textbf{PORT-US2-02} \textemdash{} \texttt{New Portfolio} 初始狀態，包含 name、base currency、\texttt{Initial holdings}、Add ETF 入口和 research-only 邊界。}{\textbf{PORT-US2-03} \textemdash{} My Portfolio \texttt{Editing} 頁面，包含 value/\allowbreak{}holdings/\allowbreak{}status、weights、share steppers、remove controls 和 \texttt{Add ETF}。}
\screenshotsSingleCompact{assets/screenshots/port-us2-04-baseline-history.png}{PORT-US2-04}{\texttt{Baseline snapshots} history、保存入口、search、兩個 saved snapshots 和所選 snapshot 的 ETF details。}
\screenshotDescSingle{\textbf{PORT-US2-04} \textemdash{} \texttt{Baseline snapshots} history、保存入口、search、兩個 saved snapshots 和所選 snapshot 的 ETF details。}
\subsection{Acceptance Criteria}
\begingroup
\small
\begin{longtable}{L{0.18\linewidth}L{0.30\linewidth}L{0.18\linewidth}L{0.12\linewidth}L{0.10\linewidth}}
\toprule
\textbf{ID} & \textbf{Requirement} & \textbf{Reference} & \textbf{Ownership} & \textbf{Priority} \\
\midrule
US2-AC01 & 首次打開或刷新時顯示載入狀態；完成後列表反映當前保存的 My Portfolio。 & US-PORT-02 & Integration & Must \\
US2-AC02 & 新建 My Portfolio 表單提供 HKD、USD、CNY、EUR 和 GBP 作為基礎貨幣。 & US-PORT-02 & Integration & Must \\
US2-AC03 & 名稱為空、存在非正數量的初始持倉或創建請求進行中時，右上角創建勾選按鈕不可用。 & US-PORT-02 & Integration & Must \\
US2-AC04 & ETF 選擇器支持代碼/\allowbreak{}名稱/\allowbreak{}發行方搜尋、排序指標、適用周期/\allowbreak{}維度、方向和重置；有效 ETF 以界面提供的默認值加入，已暫存 ETF 顯示 Added 且不可重複選擇。 & US-PORT-02 & Integration & Must \\
US2-AC05 & 無匹配搜尋顯示無結果狀態；清除搜尋後恢複 ETF 排行列表和結果數量。 & US-PORT-02 & Integration & Must \\
US2-AC06 & Initial holdings 支持 +、\ensuremath{-}、直接輸入和刪除；這些操作只修改創建草稿，最終有效數量反映在創建結果中。 & US-PORT-02 & Integration & Must \\
US2-AC07 & 創建後，My Portfolio 出現在列表中，顯示名稱、總價值、持倉數、資產配置、可用時的前一日回報，以及可用時基於 Baseline Snapshot 的偏離。 & US-PORT-02 & Integration & Must \\
US2-AC08 & 從列表打開新建 My Portfolio 後，名稱、基礎貨幣、初始持倉數量和價值與創建草稿一致。 & US-PORT-02 & Integration & Must \\
US2-AC09 & 創建流程清楚說明它創建的是研究型 My Portfolio，不會下單。 & US-PORT-02 & Integration & Must \\
US2-AC10 & 頭部展示所選 My Portfolio 的名稱、持倉數、總價值和類型；估值資料不完整時顯示對應的 Partial 或 Unavailable 狀態。 & US-PORT-02 & Integration & Must \\
US2-AC11 & 配置分組、ETF 數量和持倉權重與當前已保存持倉一致。 & US-PORT-02 & Integration & Must \\
US2-AC12 & 共享預測行為符合用戶故事 1，且使用該 My Portfolio 的已保存價值，而非此前打開的 Model Portfolio。 & US-PORT-02 & Integration & Must \\
US2-AC13 & 持倉行展示代碼、名稱、權重、份額、市值、成本和可用 ETF 元資料；選擇某行會打開匹配的 ETF 浮層。 & US-PORT-02 & Integration & Must \\
US2-AC14 & 保存會創建當前持倉、數量、價值和權重的 Baseline Snapshot。 & US-PORT-02 & Integration & Must \\
US2-AC15 & 顯示成功 Toast，UAT Before Edit 被加入 Baseline History，並帶有時間戳、價值和最高權重背景信息。 & US-PORT-02 & Integration & Must \\
US2-AC16 & 打開已保存的 Baseline Snapshot 詳情會顯示持倉數、總額、狀態以及保存的 ETF 數量和權重。 & US-PORT-02 & Integration & Must \\
US2-AC17 & Edit 頁面標明所選 My Portfolio，並顯示當前價值、持倉數、狀態和可編輯的持倉行。 & US-PORT-02 & Integration & Must \\
US2-AC18 & +、\ensuremath{-}、有效的正數量輸入及勾選提交會更新並持久化所選 ETF 數量；嘗試產生非正數量時不會提交該值。 & US-PORT-02 & Integration & Must \\
US2-AC19 & 提交成功後，刷新或重新打開仍顯示更新後的數量、價值、權重和配置。 & US-PORT-02 & Integration & Must \\
US2-AC20 & Add ETF 可按代碼、名稱或發行方搜尋，並可調整排序指標、適用維度和方向；結果數量和無結果狀態均可見。 & US-PORT-02 & Integration & Must \\
US2-AC21 & 已持有的 ETF 標記為 Added、被禁用，且無法通過選擇器重複添加。 & US-PORT-02 & Integration & Must \\
US2-AC22 & 選擇有效 ETF 後，選擇器關閉，該 ETF 以默認值加入 Edit 列表；重新打開詳情後仍存在。 & US-PORT-02 & Integration & Must \\
US2-AC23 & 移除 ETF 僅刪除該代碼，並刷新列表、價值、權重和配置。 & US-PORT-02 & Integration & Must \\
US2-AC24 & 數量編輯、添加和移除只影響對應 ETF；未操作的持倉保持原數量和成本信息。 & US-PORT-02 & Integration & Must \\
US2-AC25 & Edit 在沒有待提交修改時可以直接關閉，且不會額外改變已保存持倉。 & US-PORT-02 & Integration & Must \\
US2-AC26 & 每次創建僅添加一個 Baseline Snapshot，選中並載入其詳情，同時更新保存計數。 & US-PORT-02 & Integration & Must \\
US2-AC27 & 搜尋可按 Baseline Snapshot 標簽、日期或最高權重 ETF 代碼篩選，並可清除。 & US-PORT-02 & Integration & Must \\
US2-AC28 & 重命名拒絕空名稱，持久化有效的新標簽，並更新歷史和所選詳情。 & US-PORT-02 & Integration & Must \\
US2-AC29 & 刪除僅移除所選 Baseline Snapshot 並减少計數；若刪除當前選中項，頁面不再顯示其詳情，並選中剩餘有效項或顯示空狀態。 & US-PORT-02 & Integration & Must \\
US2-AC30 & UAT Before Edit 保留操作 2.4 前的持倉，而 UAT Current State 與操作 2.4 後的當前 My Portfolio 一致。 & US-PORT-02 & Integration & Must \\
US2-AC31 & 重新打開 Baseline snapshots 後，保留的 Baseline Snapshot 及重命名結果仍存在，被刪除的臨時 Baseline Snapshot 不再出現。 & US-PORT-02 & Integration & Must \\
US2-AC32 & 無持倉 My Portfolio 可正常打開；詳情明確顯示無持倉、配置不可用和預測不可用狀態，不以零值或示例配置冒充有效結果。 & US-PORT-02 & Integration & Must \\
US2-AC33 & 移除需要點名 My Portfolio 並說明不可撤銷的破壞性確認；取消後 My Portfolio 保持不變。 & US-PORT-02 & Integration & Must \\
US2-AC34 & 確認後從列表移除該 My Portfolio，且不刪除其他 My Portfolio。 & US-PORT-02 & Integration & Must \\
\bottomrule
\end{longtable}
\endgroup
\prdhr
\hypertarget{us-port-03}{}
\section{User Story 3 - 生成 AI Portfolio Recommendation}
\textbf{Story ID:} \texttt{US-PORT-03}
\textbf{用戶故事：} 作為用戶，我希望回答適配度問卷，獲得可解釋的 ETF 建議，檢查生成的結果並保存到 My portfolios。
\textbf{Video:} \href{https://jjpvro70sief.jp.larksuite.com/wiki/NSZ1wCHloiD8QAkYSDsj9CNtpje}{US-PORT-03 驗收視頻}
\subsection{Walkthrough}
\begin{enumerate}
\item 從 \texttt{Portfolio \textgreater{} Recommended} 打開 \texttt{Find Your Ideal Portfolio}，完成四題問卷並查看回答進度。
\item 提交答案，觀察 Recommendation 的生成進度和 pending states。
\item 查看生成結果的風險畫像、projection、portfolio metrics、Recommendation build 和 ETF holdings。
\item 返回問卷調整一個答案，重新生成並確認結果隨新答案更新。
\item 選擇 \texttt{Add to My Portfolio}，審閱配置，並選擇 \texttt{Add as new} 或 \texttt{Add to existing}。
\item 從 \texttt{My portfolios} 打開目標 portfolio，確認 Recommendation 已按所選方式保存。
\end{enumerate}
\screenshotsTripleCompact{assets/screenshots/port-us3-extra-quiz-start.png}{Recommendation quiz initial state}{\texttt{Find my portfolio} 問卷起始狀態，展示四步結構、\texttt{0/\allowbreak{}4 answered}、Goal options 和尚未啟用的 \texttt{Next}。}{assets/screenshots/port-us3-01-recommendation-quiz.png}{PORT-US3-01}{問卷進入 Question 2 後的進度狀態，顯示已完成 Goal、當前 Time step、\texttt{1/\allowbreak{}4 answered} 和 horizon options。}{assets/screenshots/port-us3-02-generation-progress.png}{PORT-US3-02}{Recommendation 正在生成的 Step 4，顯示 progress、pending header/\allowbreak{}holdings、loading projection 和 pending metrics。}
\screenshotDescsThree{\textbf{Recommendation quiz initial state} \textemdash{} \texttt{Find my portfolio} 問卷起始狀態，展示四步結構、\texttt{0/\allowbreak{}4 answered}、Goal options 和尚未啟用的 \texttt{Next}。}{\textbf{PORT-US3-01} \textemdash{} 問卷進入 Question 2 後的進度狀態，顯示已完成 Goal、當前 Time step、\texttt{1/\allowbreak{}4 answered} 和 horizon options。}{\textbf{PORT-US3-02} \textemdash{} Recommendation 正在生成的 Step 4，顯示 progress、pending header/\allowbreak{}holdings、loading projection 和 pending metrics。}
\screenshotsTripleCompact{assets/screenshots/port-us3-03-recommendation-result.png}{PORT-US3-03}{生成完成的 \texttt{Growth Focus ETF Portfolio}，包含 AI-assisted/\allowbreak{}risk/\allowbreak{}currency/\allowbreak{}holdings 頭部、projection 和 \texttt{Add to My Portfolio}。}{assets/screenshots/port-us3-extra-recommendation-build.png}{Recommendation metrics and build}{Recommendation 的 expected return、volatility、Sharpe、asset mix、generated weights 和 \texttt{Recommendation build}。}{assets/screenshots/port-us3-extra-recommendation-build-holdings.png}{Recommendation build holdings table}{\texttt{How this was built} 的 screening process、holdings table、portfolio characteristics 和 constraint note。}
\screenshotDescsThree{\textbf{PORT-US3-03} \textemdash{} 生成完成的 \texttt{Growth Focus ETF Portfolio}，包含 AI-assisted/\allowbreak{}risk/\allowbreak{}currency/\allowbreak{}holdings 頭部、projection 和 \texttt{Add to My Portfolio}。}{\textbf{Recommendation metrics and build} \textemdash{} Recommendation 的 expected return、volatility、Sharpe、asset mix、generated weights 和 \texttt{Recommendation build}。}{\textbf{Recommendation build holdings table} \textemdash{} \texttt{How this was built} 的 screening process、holdings table、portfolio characteristics 和 constraint note。}
\screenshotsTripleCompact{assets/screenshots/port-us3-extra-recommendation-etf-holdings.png}{Recommendation ETF holdings}{生成結果的完整 \texttt{ETF holdings}、target weights、deterministic-engine disclosure 和 \texttt{Back to adjust answers}。}{assets/screenshots/port-us3-04-save-recommendation.png}{PORT-US3-04}{Recommendation 的 \texttt{Add to existing} 目標選擇頁，展示可選 My Portfolio、value、holdings count、asset mix 和 drift。}{assets/screenshots/port-us3-extra-add-to-portfolio-editor.png}{Recommendation Add to My Portfolio editor}{Recommendation 進入 \texttt{Add to My Portfolio} 後的 Visual editor，下一步可選擇 \texttt{Add to existing} 或 \texttt{Add as new}。}
\screenshotDescsThree{\textbf{Recommendation ETF holdings} \textemdash{} 生成結果的完整 \texttt{ETF holdings}、target weights、deterministic-engine disclosure 和 \texttt{Back to adjust answers}。}{\textbf{PORT-US3-04} \textemdash{} Recommendation 的 \texttt{Add to existing} 目標選擇頁，展示可選 My Portfolio、value、holdings count、asset mix 和 drift。}{\textbf{Recommendation Add to My Portfolio editor} \textemdash{} Recommendation 進入 \texttt{Add to My Portfolio} 後的 Visual editor，下一步可選擇 \texttt{Add to existing} 或 \texttt{Add as new}。}
\subsection{Acceptance Criteria}
\begingroup
\small
\begin{longtable}{L{0.18\linewidth}L{0.30\linewidth}L{0.18\linewidth}L{0.12\linewidth}L{0.10\linewidth}}
\toprule
\textbf{ID} & \textbf{Requirement} & \textbf{Reference} & \textbf{Ownership} & \textbf{Priority} \\
\midrule
US3-AC01 & 問卷恰好顯示四個進度步驟及當前問題/\allowbreak{}已回答數量。 & US-PORT-03 & Integration & Must \\
US3-AC02 & 每個問題均展示輔助文本和可選擇的選項卡，含標簽、詳情、徽標及風險/\allowbreak{}強度背景。 & US-PORT-03 & Integration & Must \\
US3-AC03 & 在當前問題已有答案前，Next/\allowbreak{}See recommendation 始終禁用。 & US-PORT-03 & Integration & Must \\
US3-AC04 & 每個問題僅可選擇一個選項，且 Back/\allowbreak{}Next 導航保留所有先前答案。 & US-PORT-03 & Integration & Must \\
US3-AC05 & 完成問卷後，結果頁顯示由答案組合得出的風險等級\textemdash{}\textemdash{}Conservative、Moderate、Growth 或 Opportunistic\textemdash{}\textemdash{}並進入該等級對應的建議生成狀態。 & US-PORT-03 & Integration & Must \\
US3-AC06 & 頁面立即進入可見的生成狀態；不得短暫顯示上一次運行留下的過期建議。 & US-PORT-03 & Integration & Must \\
US3-AC07 & 進度傳達四個階段：Screen ETF universe、Optimize allocation、Draft AI explanation、Check wording。 & US-PORT-03 & Integration & Must \\
US3-AC08 & 待定的頭部、預測、構建、說明和 ETF 行均清楚標記為生成中/\allowbreak{}待定，而非填入虛構的最終數值。 & US-PORT-03 & Integration & Must \\
US3-AC09 & 在存在明確 Recommendation 結果前，保存操作不可用。 & US-PORT-03 & Integration & Must \\
US3-AC10 & 頭部展示生成結果的標題、風險等級、基礎貨幣、持倉數和 AI 輔助/\allowbreak{}僅供研究摘要。 & US-PORT-03 & Integration & Must \\
US3-AC11 & 預期年化回報、波動率、夏普比率、資產類別配置、ETF 數量和生成權重均來自返回的建議。 & US-PORT-03 & Integration & Must \\
US3-AC12 & 預測基於生成的 ETF 權重和建模金額計算，並滿足用戶故事 1 中確立的交互標準。 & US-PORT-03 & Integration & Must \\
US3-AC13 & Recommendation Build 展示本次構建的風險畫像、篩選與優化摘要、截至日期、候選範圍、刪除原因和警告，並與最終 ETF 持倉對應。 & US-PORT-03 & Integration & Must \\
US3-AC14 & AI 生成時，說明標記為 AI-generated，並能渲染段落、強調、項目符號、引用和表格，且不顯示原始 Markdown 標記。 & US-PORT-03 & Integration & Must \\
US3-AC15 & 說明內容與本次 Recommendation Build 和生成持倉一致，並清楚說明歷史資料、研究用途和非收益保證邊界。 & US-PORT-03 & Integration & Must \\
US3-AC16 & ETF 行展示代碼、名稱、資產類別、生成的目標權重和進度；顯示權重在舍入容差內合計 100\%。 & US-PORT-03 & Integration & Must \\
US3-AC17 & 返回問卷保留當前答案。重新生成會清除先前建議，並在新結果前顯示生成狀態。 & US-PORT-03 & Integration & Must \\
US3-AC18 & Recommendation 全程標明 AI 輔助和僅供研究；在用戶明確選擇 Add to My Portfolio 並完成確認前，不會保存到 My portfolios。 & US-PORT-03 & Integration & Must \\
US3-AC19 & 添加頁面使用生成的建議、建模金額、基礎貨幣和可用的最新 ETF 價格。 & US-PORT-03 & Integration & Must \\
US3-AC20 & Visual/\allowbreak{}List 編輯以及兩種目標選項均符合操作 1.5\textendash{}1.7 的標準。 & US-PORT-03 & Integration & Must \\
US3-AC21 & 使用 Add as new 時顯示確認，且新 My Portfolio 以預期的名稱、基礎貨幣和持倉數出現在 My portfolios 中。 & US-PORT-03 & Integration & Must \\
US3-AC22 & 使用 Add to existing 會合並數量：新代碼被添加，已有代碼獲得相加後的數量和加權成本。 & US-PORT-03 & Integration & Must \\
US3-AC23 & 僅打開添加頁或返回 Recommendation 不會寫入 My portfolios；只有完成目標選擇和確認後才產生對應保存結果。 & US-PORT-03 & Integration & Must \\
\bottomrule
\end{longtable}
\endgroup
\prdhr
\hypertarget{us-port-04}{}
\section{User Story 4 - 分析 Portfolio 並審閱 AI 建議}
\textbf{Story ID:} \texttt{US-PORT-04}
\textbf{用戶故事：} 作為用戶，我希望分析 My Portfolio 或 Model Portfolio 的工作草稿，選擇重要的診斷問題，並審閱 AI 輔助的考慮因素和量化 \texttt{Suggested actions}；當來源是 My Portfolio 時，我還可以編輯並應用這些操作。
\textbf{Video:} \href{https://jjpvro70sief.jp.larksuite.com/wiki/Ey4Bw3xDxiPeETk3yBBjhxOHpdf}{US-PORT-04 驗收視頻}
\subsection{Walkthrough}
\begin{enumerate}
\item 從一個 My Portfolio 選擇 \texttt{Start analysis}，查看並按需要調整草稿持倉。
\item 選擇 \texttt{Analyze draft portfolio}，查看 Analysis overview、Holdings、Exposure breakdown 和 Data coverage。
\item 查看並選擇希望建議覆蓋的 Diagnosis，然後生成 recommendations。
\item 審閱所選診斷範圍、AI considerations、candidate pool 和量化 \texttt{Suggested actions}。
\item 對 My Portfolio 來源審閱或調整目標份額，應用方案並確認 portfolio 更新。
\item 從 Model Portfolio 再次啟動 Analysis Workflow，確認 \texttt{Suggested actions} 為只讀結果。
\end{enumerate}
\screenshotsTripleCompact{assets/screenshots/port-us4-01-analysis-holdings.png}{PORT-US4-01}{\texttt{Draft Portfolio \textgreater{} Holdings} 的 analysis 前可編輯草稿，包含 shares 輸入與步進控件、刪除操作、\texttt{Add ETF}、尚未解鎖的 \texttt{Analysis /\allowbreak{} Suggestions}，以及 \texttt{Analyze draft portfolio}。}{assets/screenshots/port-us4-02-analysis-overview.png}{PORT-US4-02}{\texttt{Draft Portfolio \textgreater{} Analysis}，包含 \texttt{Data coverage}、\texttt{Analysis overview}、score、value/\allowbreak{}positions、HHI、Effective N 和多維 diagnostics。}{assets/screenshots/port-us4-extra-analysis-holdings-exposure.png}{Analysis holdings table and exposure breakdown}{Analysis 結果中的 read-only Holdings table 和 \texttt{Exposure breakdown}；該表不是運行分析前的 editable Holdings 草稿。}
\screenshotDescsThree{\textbf{PORT-US4-01} \textemdash{} \texttt{Draft Portfolio \textgreater{} Holdings} 的 analysis 前可編輯草稿，包含 shares 輸入與步進控件、刪除操作、\texttt{Add ETF}、尚未解鎖的 \texttt{Analysis /\allowbreak{} Suggestions}，以及 \texttt{Analyze draft portfolio}。}{\textbf{PORT-US4-02} \textemdash{} \texttt{Draft Portfolio \textgreater{} Analysis}，包含 \texttt{Data coverage}、\texttt{Analysis overview}、score、value/\allowbreak{}positions、HHI、Effective N 和多維 diagnostics。}{\textbf{Analysis holdings table and exposure breakdown} \textemdash{} Analysis 結果中的 read-only Holdings table 和 \texttt{Exposure breakdown}；該表不是運行分析前的 editable Holdings 草稿。}
\screenshotsDoubleCompact{assets/screenshots/port-us4-03-diagnosis-selection.png}{PORT-US4-03}{\texttt{Exposure breakdown} 的 Sector/\allowbreak{}Region/\allowbreak{}Currency，以及 \texttt{Diagnosis} 的 Found/\allowbreak{}Selected/\allowbreak{}Hidden counts、selected finding 和 generate action。}{assets/screenshots/port-us4-04-suggested-actions.png}{PORT-US4-04}{\texttt{Suggestions} 中的 selected diagnosis scope、\texttt{AI drafted} boundary、research consideration、candidate pool 和量化 \texttt{Suggested actions} preview。}
\screenshotDescsTwo{\textbf{PORT-US4-03} \textemdash{} \texttt{Exposure breakdown} 的 Sector/\allowbreak{}Region/\allowbreak{}Currency，以及 \texttt{Diagnosis} 的 Found/\allowbreak{}Selected/\allowbreak{}Hidden counts、selected finding 和 generate action。}{\textbf{PORT-US4-04} \textemdash{} \texttt{Suggestions} 中的 selected diagnosis scope、\texttt{AI drafted} boundary、research consideration、candidate pool 和量化 \texttt{Suggested actions} preview。}
\subsection{Acceptance Criteria}
\begingroup
\small
\begin{longtable}{L{0.18\linewidth}L{0.30\linewidth}L{0.18\linewidth}L{0.12\linewidth}L{0.10\linewidth}}
\toprule
\textbf{ID} & \textbf{Requirement} & \textbf{Reference} & \textbf{Ownership} & \textbf{Priority} \\
\midrule
US4-AC01 & Analysis Workspace 在 Holdings 中打開，標明來源 My Portfolio，並預載入其已保存的 ETF 持倉、數量、成本、價值和權重。 & US-PORT-04 & Integration & Must \\
US4-AC02 & 在存在分析前，Analysis 保持鎖定；在存在建議前，Suggestions 保持鎖定。 & US-PORT-04 & Integration & Must \\
US4-AC03 & +/\allowbreak{}\ensuremath{-} 會立即更新草稿；直接輸入僅在數值變化時顯示勾選確認。確認後重算價值和權重、清除舊分析/\allowbreak{}建議，且不修改已保存的 My Portfolio。 & US-PORT-04 & Integration & Must \\
US4-AC04 & Add ETF 支持目錄搜尋和排序控件；重複項標記為 Added 且無法再次添加，有效的新 ETF 可加入當前草稿。 & US-PORT-04 & Integration & Must \\
US4-AC05 & 移除 ETF 或 Clear all 僅影響工作草稿，直至顯式執行 Apply to My Portfolio。 & US-PORT-04 & Integration & Must \\
US4-AC06 & 草稿為空時 \texttt{Analyze draft portfolio} 被禁用；分析運行期間只顯示進度，不提供第二個分析操作。 & US-PORT-04 & Integration & Must \\
US4-AC07 & \texttt{Analyze draft portfolio} 切換至 \texttt{Analysis} 標簽，並持續顯示進度，直至當前草稿的分析完成或失敗。 & US-PORT-04 & Integration & Must \\
US4-AC08 & 完成後的概覽展示當前草稿可用的總價值/\allowbreak{}基礎貨幣、頭寸數、HHI、Effective N、分散化、集中度、地域、外匯、費用效率和流動性信息。 & US-PORT-04 & Integration & Must \\
US4-AC09 & 不可用指標顯示為 Unavailable/\allowbreak{}\textemdash{}；若可用資料不足以形成總體評分，該評分同樣顯示為不可用，不得以零值代替。 & US-PORT-04 & Integration & Must \\
US4-AC10 & Holdings 表顯示 Symbol、QTY、COST/\allowbreak{}UNIT、MKT VALUE、WEIGHT 和 CCY，並可橫向滚動；不可用值顯示為 \textemdash{}，不顯示占位收益。 & US-PORT-04 & Integration & Must \\
US4-AC11 & \texttt{Exposure breakdown} 顯示 Sector、Region 和 Currency，並使用同一資料在 \texttt{Pie} 和 \texttt{List} 間一致切換。 & US-PORT-04 & Integration & Must \\
US4-AC12 & 分析免責聲明始終可見；響應包含部分/\allowbreak{}不可用指標或 warning 時，頁面顯示藍色 Data coverage，說明受影響的資料範圍，而不將其稱為 AI fallback。 & US-PORT-04 & Integration & Must \\
US4-AC13 & 當前結果只對應本次草稿；後續返回 Holdings 修改數量時，舊 Analysis 和 Suggestions 會被清除並重新鎖定。 & US-PORT-04 & Integration & Must \\
US4-AC14 & 診斷項展示嚴重性、標題、指標/\allowbreak{}閾值和選擇狀態，返回項默認全選。 & US-PORT-04 & Integration & Must \\
US4-AC15 & 當前分析返回的診斷項數量、嚴重程度和選中狀態與頁面計數一致。 & US-PORT-04 & Integration & Must \\
US4-AC16 & 切換診斷項後，Found/\allowbreak{}Selected/\allowbreak{}Hidden 計數立即更新；生成結果列出的診斷範圍僅包含提交時選中的項目。 & US-PORT-04 & Integration & Must \\
US4-AC17 & 選擇零項時，UI 指示用戶至少選擇一項，且 \texttt{Generate Recommendations from  Selected} 被禁用。 & US-PORT-04 & Integration & Must \\
US4-AC18 & 分析尚未完成時生成 recommendation 的操作不可用；建議生成期間只顯示生成進度，不提供可再次觸發的生成操作。 & US-PORT-04 & Integration & Must \\
US4-AC19 & 生成切換至 Suggestions，並顯示四個進度階段：Read selected diagnosis、Draft considerations、Apply safety guard、Prepare preview。 & US-PORT-04 & Integration & Must \\
US4-AC20 & 結果標明用於生成建議的診斷項。 & US-PORT-04 & Integration & Must \\
US4-AC21 & AI 生成的考慮因素顯示 AI drafted，並帶有研究用途和 AI 可能出錯的提示。 & US-PORT-04 & Integration & Must \\
US4-AC22 & 考慮因素只覆蓋提交時選中的診斷項；未選中的診斷不會出現在本次建議範圍和對應考慮因素中。 & US-PORT-04 & Integration & Must \\
US4-AC23 & 每項考慮因素展示觀察、理由、預期影響、置信度、相關問題和僅供研究措辭，不含訂單執行語言。 & US-PORT-04 & Integration & Must \\
US4-AC24 & \texttt{Suggested actions} 展示可行性、行數、換手率、新增候選數量/\allowbreak{}代碼及風險等級。 & US-PORT-04 & Integration & Must \\
US4-AC25 & 每個 Suggested Action 行展示代碼/\allowbreak{}名稱、Held/\allowbreak{}New 狀態、當前份額/\allowbreak{}權重、目標份額/\allowbreak{}權重、數量和價值差額，以及 Model target 比較。 & US-PORT-04 & Integration & Must \\
US4-AC26 & \texttt{Suggested actions} 支持滚動查看全部 Held/\allowbreak{}New 行，並在編輯目標後立即更新對應的數量和價值差額。 & US-PORT-04 & Integration & Must \\
US4-AC27 & Recommendation candidate pool 僅展示當前建議流程允許考慮的 ETF 代碼，並與 \texttt{Suggested actions} 中的新增候選對應。 & US-PORT-04 & Integration & Must \\
US4-AC28 & 目標份額僅接受非負值，並重新計算目標權重、份額差、權重差和價值差。 & US-PORT-04 & Integration & Must \\
US4-AC29 & Reset 清除用戶覆蓋值，並恢複預覽最初給出的目標方案。 & US-PORT-04 & Integration & Must \\
US4-AC30 & Apply to My Portfolio 僅在來源為 My Portfolio、方案可行、存在可顯示操作行、價格查詢完成且所有正目標都有可用價格時啟用；不滿足條件時保持禁用。 & US-PORT-04 & Integration & Must \\
US4-AC31 & 應用會以編輯後的正目標更新已保存持倉，保留未受影響的有效持倉，關閉工作區並刷新詳情。 & US-PORT-04 & Integration & Must \\
US4-AC32 & 重新打開可證明數量和權重/\allowbreak{}價值已更新；先前分析/\allowbreak{}建議不再被視為對修改後的持倉有效。 & US-PORT-04 & Integration & Must \\
US4-AC33 & Apply to My Portfolio 請求進行中時不可重複觸發；成功後只產生一次持倉更新並返回刷新後的詳情。 & US-PORT-04 & Integration & Must \\
US4-AC34 & Model Portfolio 可以作為草稿來源完成 Holdings、Analysis、Diagnosis 和 Suggestions 的同一套 Analysis Workflow。 & US-PORT-04 & Integration & Must \\
US4-AC35 & \texttt{Suggestions} 展示本次選定的診斷範圍、AI 考慮因素和量化 \texttt{Suggested actions}，不會僅因來源為 Model Portfolio 而跳過結果。 & US-PORT-04 & Integration & Must \\
US4-AC36 & Model Portfolio 的 \texttt{Suggested actions} 以只讀列表展示，每行可查看 Current、Target、Change 和價值差額，但不提供 Target shares 輸入或步進控件。 & US-PORT-04 & Integration & Must \\
US4-AC37 & Model Portfolio 結果不顯示 Reset 或 Apply to My Portfolio；關閉工作區不會修改該 Model Portfolio，也不會創建或更新 My Portfolio。 & US-PORT-04 & Integration & Must \\
\bottomrule
\end{longtable}
\endgroup
\prdhr
\hypertarget{us-port-05}{}
\section{User Story 5 - 根據 Baseline Snapshot 審閱}
\textbf{Story ID:} \texttt{US-PORT-05}
\textbf{用戶故事：} 作為用戶，我希望將今天的 My Portfolio 與有效的歷史 Baseline Snapshot 比較，了解配置偏離，編輯基於 Baseline Snapshot 的調整方案並應用。
\textbf{Video:} \href{https://jjpvro70sief.jp.larksuite.com/wiki/H1ULwBLKeiaEIKkEbQDjhRiepue}{US-PORT-05 驗收視頻}
\subsection{Walkthrough}
\begin{enumerate}
\item 從一個 My Portfolio 打開 \texttt{Portfolio Review}，瀏覽可用的 Baseline Snapshots，並區分 \texttt{No changes} 與可審閱狀態。
\item 選擇一個有變化的 Baseline Snapshot，查看 portfolio 摘要和 holdings drift。
\item 進入 \texttt{Analysis}，閱讀 drift status、driver summary、allocation changes 和 ETF quantity changes。
\item 打開 \texttt{Suggestions}，審閱並按需要調整 Baseline-based target shares。
\item 選擇 \texttt{Apply to My Portfolio}，返回 portfolio 並確認調整結果已更新。
\end{enumerate}
\screenshotsTripleCompact{assets/screenshots/port-us5-01-baseline-selection.png}{PORT-US5-01}{\texttt{Portfolio Review \textgreater{} Baseline}，同時顯示 \texttt{No changes} snapshot、selected snapshot、ETF summary 和 \texttt{Review allocation changes}。}{assets/screenshots/port-us5-02-baseline-deviation.png}{PORT-US5-02}{所選 Baseline Snapshot 的 current/\allowbreak{}baseline value、absolute drift 和 \texttt{Holding drift} table，逐行比較 Now、Baseline 和 Drift。}{assets/screenshots/port-us5-03-review-analysis.png}{PORT-US5-03}{Baseline comparison 的 \texttt{Analysis}，包含 status、allocation drift、current/\allowbreak{}baseline values、largest move 和 driver summary。}
\screenshotDescsThree{\textbf{PORT-US5-01} \textemdash{} \texttt{Portfolio Review \textgreater{} Baseline}，同時顯示 \texttt{No changes} snapshot、selected snapshot、ETF summary 和 \texttt{Review allocation changes}。}{\textbf{PORT-US5-02} \textemdash{} 所選 Baseline Snapshot 的 current/\allowbreak{}baseline value、absolute drift 和 \texttt{Holding drift} table，逐行比較 Now、Baseline 和 Drift。}{\textbf{PORT-US5-03} \textemdash{} Baseline comparison 的 \texttt{Analysis}，包含 status、allocation drift、current/\allowbreak{}baseline values、largest move 和 driver summary。}
\screenshotsTripleCompact{assets/screenshots/port-us5-extra-allocation-movements.png}{Baseline allocation and ETF movements}{\texttt{Industry allocation changes}、\texttt{ETF movement detail} 和 quantity-difference summary。}{assets/screenshots/port-us5-extra-quantity-changes.png}{Baseline quantity changes}{\texttt{Quantity changes} 中各 ETF 的 baseline/\allowbreak{}current/\allowbreak{}target quantities 和 value differences。}{assets/screenshots/port-us5-04-adjustment-suggestions.png}{PORT-US5-04}{可編輯的 baseline-based \texttt{Recommended ETF changes}，包含 target-share controls、\texttt{Reset suggestion} 和 \texttt{Apply to My Portfolio}。}
\screenshotDescsThree{\textbf{Baseline allocation and ETF movements} \textemdash{} \texttt{Industry allocation changes}、\texttt{ETF movement detail} 和 quantity-difference summary。}{\textbf{Baseline quantity changes} \textemdash{} \texttt{Quantity changes} 中各 ETF 的 baseline/\allowbreak{}current/\allowbreak{}target quantities 和 value differences。}{\textbf{PORT-US5-04} \textemdash{} 可編輯的 baseline-based \texttt{Recommended ETF changes}，包含 target-share controls、\texttt{Reset suggestion} 和 \texttt{Apply to My Portfolio}。}
\subsection{Acceptance Criteria}
\begingroup
\small
\begin{longtable}{L{0.18\linewidth}L{0.30\linewidth}L{0.18\linewidth}L{0.12\linewidth}L{0.10\linewidth}}
\toprule
\textbf{ID} & \textbf{Requirement} & \textbf{Reference} & \textbf{Ownership} & \textbf{Priority} \\
\midrule
US5-AC01 & Portfolio Review 在 Baseline 標簽打開，且僅列出該 My Portfolio 保存的快照。 & US-PORT-05 & Integration & Must \\
US5-AC02 & UI 說明 Baseline Snapshot 是已保存的參考配置，而不是手動配置的目標。 & US-PORT-05 & Integration & Must \\
US5-AC03 & UAT Current State 明顯禁用/\allowbreak{}變暗，標記為 No changes，並說明當前持倉與 Baseline Snapshot 相同。 & US-PORT-05 & Integration & Must \\
US5-AC04 & 先選中 UAT Before Edit 再點擊無變化的 Baseline Snapshot 時，UAT Current State 不會被選中，原選擇保持不變。 & US-PORT-05 & Integration & Must \\
US5-AC05 & 標記為 No changes 的 Baseline Snapshot 不會提供進入 Analysis 的操作，且頁面說明其持倉與當前狀態一致。 & US-PORT-05 & Integration & Must \\
US5-AC06 & 即使 Baseline Snapshot 無法選中以供審閱，Show ETF details 仍可用於檢查。 & US-PORT-05 & Integration & Must \\
US5-AC07 & 有變化的 Baseline Snapshot 變為 Selected，並獲得選中樣式。 & US-PORT-05 & Integration & Must \\
US5-AC08 & 展開的詳情展示保存的 ETF 數量和權重；折疊僅隱藏詳情面板。 & US-PORT-05 & Integration & Must \\
US5-AC09 & 偏離結果同時列出新增和已移除的 ETF，不得因某代碼只存在於當前持倉或 Baseline Snapshot 中而遺漏。 & US-PORT-05 & Integration & Must \\
US5-AC10 & 偏離表按重大變動排序，展示每個代碼的 Baseline Snapshot 權重、當前權重、數量差和帶符號的權重偏離。 & US-PORT-05 & Integration & Must \\
US5-AC11 & 摘要標明所選 Baseline Snapshot 的名稱、日期及總絕對偏離；Review allocation changes 現已啟用。 & US-PORT-05 & Integration & Must \\
US5-AC12 & Analysis 清楚標明正在比較當前持倉與所選 Baseline Snapshot。 & US-PORT-05 & Integration & Must \\
US5-AC13 & Analysis 顯示 Within range、Monitor drift 或 Review recommended 狀態，並展示當前價值、Baseline Snapshot 價值、最大變動和百分點偏離。 & US-PORT-05 & Integration & Must \\
US5-AC14 & \texttt{Driver summary} 標識最大增加/\allowbreak{}减少的敞口，並說明價格變動和/\allowbreak{}或數量變動是否有貢獻。 & US-PORT-05 & Integration & Must \\
US5-AC15 & \texttt{Industry allocation changes} 按 ETF 資產類別匯總 Baseline Snapshot 權重和當前權重，並展示帶符號的偏離。 & US-PORT-05 & Integration & Must \\
US5-AC16 & \texttt{ETF movement detail} 展示 Baseline Snapshot 權重、當前權重和數量，以及每項重大變動的通俗原因。 & US-PORT-05 & Integration & Must \\
US5-AC17 & \texttt{Quantity changes} 為關鍵行展示 Baseline Snapshot 數量、當前數量、目標數量、單位差和價值差。 & US-PORT-05 & Integration & Must \\
US5-AC18 & Baseline Snapshot 比較不得顯示 AI-generated、AI drafted 等 AI 標記，並明確呈現為基於所選快照的配置比較。 & US-PORT-05 & Integration & Must \\
US5-AC19 & Suggestions 僅在所選 Baseline Snapshot 的預覽載入完成後顯示目標方案；載入或切換範圍時不顯示上一份過期預覽。 & US-PORT-05 & Integration & Must \\
US5-AC20 & 範圍卡點名 Baseline Snapshot，並報告變更行數和總行數。 & US-PORT-05 & Integration & Must \\
US5-AC21 & 目標行展示當前份額/\allowbreak{}權重、Baseline Snapshot 權重、計算出的目標份額/\allowbreak{}權重、帶符號的變動、價值差和 Hold/\allowbreak{}Adjust 狀態。 & US-PORT-05 & Integration & Must \\
US5-AC22 & 默認方案按所選 Baseline Snapshot 生成可審閱的目標份額；Baseline Snapshot 中不存在的 ETF 可顯示為零目標。 & US-PORT-05 & Integration & Must \\
US5-AC23 & 目標值可編輯、非負，並立即重新計算顯示的權重和差額。 & US-PORT-05 & Integration & Must \\
US5-AC24 & Reset suggestion 移除覆蓋值，並重新設定計算出的默認目標。 & US-PORT-05 & Integration & Must \\
US5-AC25 & 在目標計算完成前 Apply to My Portfolio 被禁用；目標方案可用後才允許提交。 & US-PORT-05 & Integration & Must \\
US5-AC26 & Apply to My Portfolio 成功後關閉審閱並刷新 My Portfolio 詳情，詳情顯示已更新的保存持倉。 & US-PORT-05 & Integration & Must \\
US5-AC27 & 重新打開可證明已應用目標份額，且此前選中的 Baseline Snapshot 顯示更新後的比較結果。 & US-PORT-05 & Integration & Must \\
US5-AC28 & Apply to My Portfolio 請求進行中時不可重複觸發；完成後只產生一次持倉更新，重新打開詳情和 Portfolio Review 均顯示同一最終狀態。 & US-PORT-05 & Integration & Must \\
\bottomrule
\end{longtable}
\endgroup
